\section{Cook--Reckhow's source framework and relative comparisons}
\label{sec:p07-cook-reckhow-source}

Cook and Reckhow's 1979 paper compares representations and proof systems,
not algorithms that find shortest proofs. Its general proof-system
definition, uniform simulation and separate formula/line measures make
precise what an efficiency comparison establishes
\cite[\S\S1--5]{CookReckhow1979}. The following source treatment also
supplies its natural-deduction and extension interfaces, its pigeonhole
example, and the exact boundary between polynomial boundedness and
\(\mathsf P=\mathsf{NP}\).

\subsection{Encoded proof maps, verifiers and polynomial boundedness}

\begin{hbconvention}{def:p07-cr-encoding}{Source string lengths and finite alphabets}
The source uses finite alphabets \(\Sigma_1,\Sigma_2\) and total
polynomial-time functions \(f:\Sigma_1^*\to\Sigma_2^*\).
Write \(|x|\) for string length in this explicitly source-bound
convention. It is not set cardinality. A fixed binary encoding of the
finite alphabet changes lengths by a constant factor. Atoms are encoded
as a letter followed by a binary identifier; delimiters make formulas,
proof lines and tuples uniquely parseable. Distinct atom identifiers
remain distinct, regardless of their numerical magnitude.

A paired code \(\langle y,w\rangle\) explicitly contains both its
claimed conclusion y and certificate w, with length
\(O(|y|+|w|+1)\). On a malformed proof code use a fixed default
\(y_0\) in the target language. The computation time of a total
proof map bounds its output length. A polynomial-time relation on
\((w,y)\) does not supply an output-recovery search from w alone.
\end{hbconvention}

\begin{hbdefinition}{def:p07-cr-proof-map}{The source definition for an arbitrary language}
For a nonempty string language \(L\subseteq\Sigma_2^*\), a proof
system is a polynomial-time computable surjection
\(f:\Sigma_1^*\to L\). An f-proof of y is any x with f(x)=y.
It is polynomially bounded when a single polynomial q satisfies:
every \(y\in L\) has some proof x with \(|x|\leq q(|y|)\).
The definition permits systems not presented by a familiar calculus.
For \(L=\mathsf{TAUT}\), it gives the propositional proof systems
of Definition~\ref{def:p07-cook-reckhow}.
The empty language is handled separately: no total map from the nonempty
set of all strings can have empty range
\cite[defs.~1.2--1.3, prop.~1.4]{CookReckhow1979}.
\end{hbdefinition}

\begin{hbtheorem}{thm:p07-cr-map-relation}{proved=annex:p7.cr.map-relation}{The exact map--relation encoding correspondence}
A proof map f gives a polynomial-time relation R with
\(R(x,y)\) exactly when f(x)=y. Conversely, a polynomial-time
decidable relation R whose existential range is exactly a nonempty L
gives the total polynomial-time proof map
\[
 f_R(\langle y,w\rangle)=
 \begin{cases}y,&R(w,y),\\ y_0,&\text{otherwise},\end{cases}
\]
with the same default on malformed codes. Polynomial boundedness
corresponds to polynomially bounded witness existence for R, with y
included in the proof code. Neither this construction nor its
equivalence searches an unbounded set of possible outputs.
\end{hbtheorem}

\begin{hbtheorem}{thm:p07-cr-np-languages}{proved=annex:p7.cr.np-languages}{Source Proposition 1.4, including the empty case}
A language L belongs to \(\mathsf{NP}\) if and only if L is empty
or has a polynomially bounded source proof system. This is a
short-certificate existence theorem, not a deterministic proof-search
algorithm.
\end{hbtheorem}

\begin{hbtheorem}{thm:p07-cr-np-conp}{proved=annex:p7.cr.np-conp}{The exact propositional complexity equivalence}
\[
\begin{aligned}
 &\text{some propositional proof system is polynomially bounded}\\
 &\quad\Longleftrightarrow\quad \mathsf{TAUT}\in\mathsf{NP}
 \quad\Longleftrightarrow\quad\mathsf{NP}=\mathsf{coNP}.
\end{aligned}
\]
The final equivalence uses the NP-completeness reduction for SAT in
Theorem~\ref{thm:p04-cook-levin}, the same external reduction used
in the source's Proposition 1.1. It does not identify
\(\mathsf{NP}=\mathsf{coNP}\) with \(\mathsf P=\mathsf{NP}\).
If \(\mathsf P=\mathsf{NP}\), polynomial boundedness follows;
the converse is not established by this theorem.
\end{hbtheorem}

\begin{hbdefinition}{def:p07-cr-source-simulation}{Uniform simulation and shortest-proof comparison}
For proof maps f and h with the \emph{same} target language, h
p-simulates f when a single polynomial-time total transformation g
satisfies \(h(g(x))=f(x)\) for all strings x, including codes sent
to the default conclusion. The source's p-equivalence means simulation
in both directions. Separately, put
\(s_f(y)=\min\{|x|:f(x)=y\}\). A bound
\(s_h(y)\leq q(s_f(y))\) compares existence of short proofs;
it does not by itself give an algorithm translating a supplied proof.
Different connective alphabets generally have different literal formula
languages, so a comparison between them must state its translation or
its length-bound meaning rather than silently use this same-range definition.
\end{hbdefinition}

\begin{hbtheorem}{thm:p07-cr-simulation-properties}{proved=annex:p7.cr.simulation-properties}{Composition, length transfer and what remains nonuniform}
P-simulation is reflexive and transitive. Mutual p-simulation is an
equivalence relation. If h p-simulates f and f is polynomially bounded,
then h is polynomially bounded. The translation also gives a polynomial
shortest-proof bound from f to h. The shortest-proof inequality alone
does not specify a translating procedure or a running-time bound for
finding its witnesses.
\end{hbtheorem}

\subsection{The source's four size measures and Frege systems}

\begin{hbdefinition}{def:p07-cr-source-measures}{Formula occurrences, lines, maximum width and encoded bits}
For a formula A, \(l(A)\) counts atom and nullary-connective
\emph{occurrences}, not distinct atoms or encoded bits. For a sequence
of formulas, it sums that count. For a derivation \(\pi\),
\(\lambda(\pi)\) counts its lines and
\(\rho(\pi)=\max_i l(A_i)\) is its maximum formula count.
The separate measure \(|\pi|\) is its encoded string length.
For the defining-set bounds below, also let \(m(A)\) count all
nodes of the formula parse tree, including unary-connective occurrences.
Thus \(l(A)\leq m(A)\leq |A|\) up to fixed encoding constants;
an arbitrarily long unary chain can have l-count one.
For natural-deduction lines, l counts both the ordered hypotheses
and the conclusion. Parent indices and rule annotations also cost
encoded bits; repeated use of one earlier line does not duplicate
that line in a sequence proof.

These source measures need not be interchangeable: a long binary
atom name counts once in l; a proof can have few lines containing
enormous formulas; a DAG can share a subproof that occurs many times
in its unfolded tree. The bounds below retain their stated measures
\cite[p.~39]{CookReckhow1979}.
\end{hbdefinition}

\begin{hbdefinition}{def:p07-cr-source-frege}{Source Frege rules and implicational completeness}
Fix a finite adequate connective set \(\kappa\), with arities at
most two. A rule is a finite schematic inference
\((C_1,\ldots,C_r)/D\) whose premise formulas semantically entail
its conclusion. Instances use simultaneous formula substitution.
A source Frege system \(\mathcal F\) is a finite collection of
these sound rules that is \emph{implicationally complete}: every
valid finite-premise consequence has a derivation. Zero-premise
rules are axiom schemata. This is the hypothesis needed when replacing
one rule by a fixed derivation in another system, not merely a
theoremhood-completeness assumption.

A proof code records the formula sequence, rule choice, premise
indices and instance substitution, with no hypotheses for theoremhood.
Finite template matching and local inference checking are polynomial
in that code. Return its last formula when valid and a fixed tautology
otherwise. Soundness and implicational completeness give a source
proof map with exactly \(\mathsf{TAUT}_\kappa\) as its range.
\end{hbdefinition}

\begin{hbtheorem}{thm:p07-cr-frege-comparison}{proved=annex:p7.cr.frege-comparison}{Theorem 2.3: fixed-basis uniform Frege comparisons}
Given source Frege systems \(\mathcal F_1,\mathcal F_2\) over
the same \(\kappa\), one polynomial-time transformation takes a
derivation of B from \(A_1,\ldots,A_r\) in the first into a
derivation of the same B from the same hypotheses in the second,
with \(\lambda(\pi')\leq c\lambda(\pi)\) and
\(\rho(\pi')\leq c\rho(\pi)\). The associated proof maps
p-simulate each other, including default inputs. Their constants
depend on the two fixed rule sets, not on the individual derivation.
The different-connective-set Frege generalization mentioned in the
paper is attributed there to Reckhow's separate thesis; it is not
inferred by naive expansion of duplicated connectives.
\end{hbtheorem}

\begin{hbdefinition}{def:p07-cr-source-nd}{Source ordered-context natural deduction}
A source natural-deduction line is \(\Gamma\to A\), where
\(\Gamma=(A_1,\ldots,A_r)\) is a finite \emph{sequence}.
It represents the left-associated disjunction
\(L^*=\bigvee(\neg A_1,\ldots,\neg A_r,A)\), or its
right-associated form \(L^\#\). With no hypotheses both are A.
A source rule is a finite valid inference of such lines; an instance
may simultaneously substitute formulas and prefix the \emph{same}
context sequence \(\Delta\) to every premise and conclusion.
A system \(\mathcal N\) has finitely many such rules and is
implicationally complete for line consequences. Proofs are sequences
of reusable lines, not obligatorily unfolded trees
\cite[\S3]{CookReckhow1979}.
\end{hbdefinition}

\begin{hbtheorem}{thm:p07-cr-nd-comparison}{proved=annex:p7.cr.nd-comparison}{Theorem 3.3 and the source Frege--ND bridge}
Two source natural-deduction systems over the same adequate basis
have uniform conversions with \(\lambda'\leq c\lambda\) and
\(\rho'\leq c\rho\). Every such system and every source
Frege system on that basis p-simulate one another for theoremhood.
The construction preserves ordered contexts and prefixes; it does
not eliminate nested deductions by repeatedly duplicating an entire
earlier proof. The conclusion is not a tree-proof lower-bound theorem.
\end{hbtheorem}

\subsection{Definitional extension and the source compression arguments}

\begin{hbdefinition}{def:p07-cr-source-extension}{Exact extension freshness and final-formula conditions}
The source extension \(e\mathcal F\) may add
\(P\leftrightarrow C\) when P is absent from C, every earlier
line, every hypothesis, and the final conclusion. Later intermediate
lines may use P. If biconditional is not primitive, use its fixed
representation over \(\kappa\). Definitions are introduced in
dependency order, so they are acyclic. This is the source interface
of the repaired \hyperref[sheet:extended-frege-cpl]{extended-Frege sheet},
not permission to conclude an arbitrary newly declared atom.
\end{hbdefinition}

\begin{hbtheorem}{thm:p07-cr-extension-elimination}{proved=annex:p7.cr.extension-elimination}{Soundness and source Proposition 4.3}
Every extension derivation is sound for its original hypotheses and
conclusion. If it has m definitions, there is an ordinary Frege
derivation with at most \(\lambda(\pi)+cm\) lines of that same
conclusion from those same hypotheses. Definition elimination can
expand its formula lengths exponentially, so this line-count claim
is not a p-simulation of extended Frege by ordinary Frege.
\end{hbtheorem}

\begin{hbdefinition}{def:p07-cr-gate-definitions}{The source defining-set construction}
Associate to each nonatomic formula B a distinct proxy atom \(P_B\),
and set \(P_Q=Q\) for an atom Q. Ordinary atoms and newly introduced
proxies occupy disjoint name spaces. For connective sets
\(\kappa,\kappa'\), fix a \(\kappa'\)-formula computing the
truth function of each \(\kappa\) connective. The set
\(\operatorname{def}_{\kappa'}(A)\) has one defining formula
per distinct nonatomic subformula: a nullary proxy equals the fixed
constant representation; a unary proxy equals its connective's
representation on its child proxy; and a binary proxy equals its
representation on its two child proxies. This uses child atoms,
not recursively expanded child formulas.

These definitions uniquely determine all proxies from the original
atom values and imply \(P_A\leftrightarrow A\) when both sides
are expressed in one basis. Their total atom-occurrence count is
\(O(m(A))\), including gates along unary chains. In an encoded
transformation assign proxies fresh
integer identifiers by traversing the actual subformula records;
their names need not contain an entire formula as a literal string.
\end{hbdefinition}

\begin{hbtheorem}{thm:p07-cr-gate-compression}{proved=annex:p7.cr.gate-compression}{Lemmas 4.10--4.11 with the actual introduction order}
From a derivation of \(P_A\) under
\(\operatorname{def}_{\kappa}(A)\) in \(e\mathcal F\), the
source skeleton construction produces a derivation of the same
\(P_A\) under \(\operatorname{def}_{\kappa'}(A)\) in
\(e\mathcal F'\), with \(\lambda'\leq c\lambda\) and
constant \(\rho'\leq d\). The hypothesis, ordinary-rule and
extension-rule cases each use finitely many fixed gate templates.
Unneeded gates inside a large substitution formula are not expanded.
Extension gates are introduced online in dependency order, not
preloaded before a source fresh variable is introduced.

Conversely, replacing each hypothesis proxy for a subformula E of
A by E converts a width-r derivation of \(P_A\) into a proof of
A with at most \(\lambda+c m(A)\) lines and maximum formula
count at most \((r+c)l(A)\). Fresh auxiliary definitions are first
renamed away from the substituted formulas.
\end{hbtheorem}

\begin{hbtheorem}{thm:p07-cr-ef-line-compression}{proved=annex:p7.cr.ef-line-compression}{Short lines with every formula constructor counted}
For every proof \(\pi\) of A in a fixed extended Frege system,
there is a proof of that same A with
\[
 \lambda(\pi')\leq c(\lambda(\pi)+m(A)),\qquad
 \rho(\pi')\leq c l(A).
\]
This transformation is effective and polynomial in the encoded input
proof. It does not assert that an unextended Frege proof admits the
same compression.
\end{hbtheorem}

\begin{hbtheorem}{thm:p07-cr-ef-relative-bounds}{proved=annex:p7.cr.ef-relative-bounds}{Relative extension bounds and the exact complexity consequence}
Let L be a nondecreasing natural-number bound with \(L(n)\geq n\).
If every tautology A over \(\kappa\) has an \(e\mathcal F\)-proof
with at most \(L(m(A))\) lines, then every tautology B over
\(\kappa'\) has an \(e\mathcal F'\)-proof with
\[
 \lambda(\pi')\leq cL(c m(B)),\qquad \rho(\pi')\leq c l(B).
\]
An arbitrary length bound can first be replaced by its nondecreasing
prefix maximum; polynomial bounds remain polynomial under that operation.
Consequently one extended Frege system is polynomially bounded exactly
when every such fixed-basis system is, and polynomial line boundedness
is equivalent to polynomial encoded-proof boundedness for these systems.
If \(\mathsf{NP}\ne\mathsf{coNP}\), none has polynomially bounded
line counts for all tautologies; neither do Frege or the source ND
systems. This is the corrected premise of source Corollary 4.7
\cite{CookReckhow1979,CookReckhow1979AuthorCorrection}.
It is a short-proof-existence comparison across different formula
languages, not an unstated same-range p-simulation or proof-search algorithm.
\end{hbtheorem}

\begin{hbdefinition}{def:p07-cr-source-substitution}{Substitution Frege is theorem-only}
The source system \(s\mathcal F\) adds simultaneous substitution
of formulas for atoms to a Frege system, but permits only proofs with
no hypotheses. Substitution preserves tautologies; it does not preserve
a consequence from unchanged hypotheses. For example P does not entail
Q merely because substituting Q for P changes P to Q.
\end{hbdefinition}

\begin{hbtheorem}{thm:p07-cr-substitution-simulation}{proved=annex:p7.cr.substitution-simulation}{Theorem 5.3: a uniform extension-to-substitution conversion}
The source \(s\mathcal F\) p-simulates \(e\mathcal F\) on the
same formula language. The resulting proof has
\(\lambda'\leq c\lambda\rho\) and
\(\rho'\leq c\lambda\rho\). Its elimination substitutions
are applied in reverse definition order, so they do not recursively
expand the entire definition DAG. No converse simulation is asserted.
\end{hbtheorem}

\subsection{The pigeonhole example: clause inputs, sequences and expanded trees}

\begin{hbdefinition}{def:p07-cr-source-php}{The paper's indexed pigeonhole family}
For \(n\geq2\), use \(P_{ij}\) with \(1\leq i\leq n\),
\(1\leq j\leq n-1\). The source clause family
\(\mathcal P_n\) consists of each row disjunction
\(P_{i1}\lor\cdots\lor P_{i,n-1}\), and the collision clauses
\(\neg P_{ij}\lor\neg P_{kj}\) for \(i<k\) and each hole j.
It does not include functionality clauses forbidding one pigeon
from having two true hole variables. Its tautology is
\(A_n=\neg\bigwedge\mathcal P_n\).
After renaming atoms, this is the \(n-1\)-hole instance of
Definition~\ref{def:p07-pigeonhole-cnf}; the one-index shift is explicit.
\end{hbdefinition}

\begin{hbtheorem}{thm:p07-cr-php-compression}{proved=annex:p7.cr.php-compression}{The source pigeonhole induction with and without abbreviations}
From \(\mathcal P_n\), substitute
\(B_{ij}=P_{ij}\lor(P_{i,n-1}\land P_{nj})\),
\(i<n,j<n-1\), for the variables of \(\mathcal P_{n-1}\).
The resulting smaller pigeonhole clauses follow by explicit finite
Boolean inferences. Iterating and using the n=2 contradiction gives
Frege proofs with \(O(n^4)\) lines, whose repeated-substitution
formulas can have exponentially many leaf occurrences. With fresh
acyclic proxy atoms for each induction stage, the derivation from
the individual pigeonhole clauses has \(O(n^4)\) lines and
\(O(n)\) maximum leaf count. A guarded proof and a balanced defining
set discharge those inputs, producing a proof of \(A_n\) in
\(O(n^5\log(n+2))\) encoded bits. Its final formula itself has
\(\Theta(n^3)\) leaf occurrences; the smaller width bound applies
to the derivation, not to every line of the full proof. These are
bounds for the construction, not a
lower bound on the shortest unextended Frege proofs.
\end{hbtheorem}

\begin{hbtheorem}{thm:p07-cr-small-php}{proved=annex:p7.cr.small-php}{A complete small resolution DAG and its unfolded tree}
The three-pigeon/two-hole instance \(\mathcal P_3\) has the
19-node resolution DAG in its reciprocal proof: nine input clauses
and ten derived clauses, ending in the empty clause. Its binary-tree
unfolding has 21 clause occurrences because two input nodes are
reused. These are sizes of the displayed derivation and its unfolding,
not claims about either system's minimum proof size or an exponential
separation between tree-like and DAG resolution.
\end{hbtheorem}

\paragraph{Lower bounds and dated open status.}
The separately attributed Haken result remains
Theorem~\ref{thm:p07-haken-resolution-lower-bound} with its
complete counting proof in Annex~\ref{annex:p7.complexity.1}. Its exponential resolution bound does
not imply a Frege or extended-Frege lower bound, and a lower bound for
one fixed system does not prove \(\mathsf{NP}\ne\mathsf{coNP}\).
The current edition's dated unrestricted-Frege/extended-Frege status
is recorded separately in
Theorem~\ref{thm:p07-frege-lower-bound-boundary}; it is not a result
of the 1979 paper. The original paper's conjecture that no system
is polynomially bounded must likewise remain a conjecture, not a
proof of a complexity-class separation. Haken 1985, the Beame--Pitassi
survey and Kraj\'i\v{c}ek's monograph retain their separate source
allocations; none is represented as read or fully reconstructed here.
